\section{Related Work}

In this section, we discuss other techniques to train models at scale.

\paragraph{Parallelism for Large Models.} Pipeline model parallelism is a common technique used to train large models. Pipeline parallelism comes in a few flavors: the mode discussed in this paper uses flushes to ensure \emph{strict} optimizer semantics. TeraPipe~\cite{li2021terapipe} exposes fine-grained pipeline parallelism across tokens in a single training sequence for auto-regressive models like GPT. PipeTransformer~\cite{he2021pipetransformer} elastically adjusts the degree of pipelining and data parallelism by freezing layers with ``stable'' weights, and instead dedicates resources to train the remaining ``active'' layers. HetPipe~\cite{park2020hetpipe} uses a combination of pipeline and data parallelism on a set of heterogeneous accelerators. Pipeline parallelism can also be implemented with relaxed semantics: PipeDream-\texttt{2BW}~\cite{narayanan2021memory} maintains two weight versions and guarantees 1-stale weight updates without expensive flushes, while PipeMare~\cite{yang2021pipemare} and Kosson et al.~\cite{pb2021kosson} use asynchoronous pipeline parallelism. These techniques have improved throughput compared to the techniques with pipeline flushes considered in this paper, but potentially at the cost of convergence rate or final accuracy. Moreover, pipeline parallelism in isolation can still only scale to a number of devices equal to the number of layers in the model, which is limiting for certain model architectures.

PipeDream~\cite{narayanan2019pipedream} combined pipeline parallelism and data parallelism in a principled way to reduce cross-device communication. DeepSpeed~\cite{deepspeed3Dparallelism} combined pipeline parallelism with tensor and data parallelism to train models with up to a trillion parameters, but with lower throughput than what was shown in this paper (52\% vs. 36\% of peak) for a few reasons: operator fusion to keep most of the operator graph compute-bound, a more-efficient pipeline parallelism schedule to minimize the pipeline bubble size, fast hardware (A100 vs. V100 GPUs and high-bandwidth links between GPUs on the same and different servers), and scaling to more GPUs. We want to emphasize that this higher throughput makes estimated training times much more practical (about 3 months); an aggregate throughput of 37.6 petaFLOP/s would take about 40 months to train an equivalently-sized model. We can scale to larger models as well, but would need more GPUs to keep training time practical.

Mesh-TensorFlow~\cite{mesh_tf} proposes a language for easily specifying parallelization strategies that combine data and model parallelism. Switch Transformers~\cite{fedus2021switch} used Mesh-Tensorflow to train a sparsely activated expert-based model with 1.6 trillion parameters, with improved pre-training speed over the T5-11B model~\cite{T5}.

\paragraph{Sharded Data Parallelism.} As part of performance optimizations for MLPerf 0.6~\cite{mattson2019mlperf}, sharded data parallelism~\cite{kumar2019scale, xu2020automatic}, where optimizer state is sharded over data-parallel workers, was introduced. This method has two advantages: (a) it does not introduce extra communication over vanilla data parallelism, and (b) it divides the optimizer's computation and memory cost across the data-parallel partitions. ZeRO~\cite{rajbhandari2019zero,rajbhandari2021zero} extends this idea: weight parameters and gradients are sharded across data-parallel workers as well, and workers fetch relevant state from their ``owning'' workers before performing computations. This adds additional communication, which can be partially hidden by carefully overlapping computation and communication. However, this can become harder if tensor parallelism is not used or the batch size is not large enough to hide the extra communication overhead (Figure~\ref{fig:zero3_comparison_throughputs}). ZeRO-Infinity~\cite{rajbhandari2021zero} uses NVMe to efficiently swap parameters, enabling the training of very large models on a small number of GPUs. We note that using a small number of GPUs for training a very large model results in unrealistic training times (e.g., thousands of years to converge).


\paragraph{Automatic Partitioning.} FlexFlow~\cite{flexflow}, PipeDream~\cite{narayanan2019pipedream}, DAPPLE~\cite{fan2021dapple}, and Tarnawski et al.~\cite{tarnawski2020efficient} all auto-partition model training graphs over multiple devices with the help of cost models. However, each of these do not consider \emph{all} the parallelism dimensions considered in this paper: pipeline and tensor model parallelism, data parallelism, microbatch size, and the effect of memory-savings optimizations like activation recomputation on the training of models larger than the memory capacity of an accelerator. These added dimensions increase the search space that needs to be explored. Gholami et al.~\cite{gholami2018integrated} show how communication costs for combinations of data and model parallelism can be modeled.

\paragraph{HPC for Model Training.} Goyal et al.~\cite{goyal2017accurate} and You et al.~\cite{you2018imagenet} both demonstrate the use of High Performance Computing techniques to train highly-accurate ImageNet models in minutes. However, the image classification models considered fit comfortably on a single accelerator, rendering model parallelism unnecessary, support very large batch sizes ($>32$k) that allow scaling data parallelism to large worker counts with infrequent communication, and are composed of compact convolutional layers that are inherently amenable to data-parallel communication.
